% META: {"id":"1NSI-TC-CO-035","chapitre":"1NSI-TYPES-CONSTRUITS","type_objet":"corrige","capacites":["1NSI-TYPES-CONSTRUITS-C4"],"mode_creation":"ex_nihilo","sources_inspiration":[],"status":"approved","exercice_ref":"1NSI-TC-EX-035","origin":"needs_review","status_history":["needs_review","audited","accepted","approved"],"release_acceptance":"RELEASE_OWNER_BATCH_ACCEPTANCE","acceptance_closure_digest":"sha256:9b3ccf9a81c5520fbb7e03b7b2d3e7bf2057a02834b6d908bf3be5f49f3b553f","acceptance_receipt_digest":"sha256:4fad86f0282e0fa4e0419cca1ad6379ae7af5a280c04a4a90880f90a1c044242"}
% BEGIN-VERIFY
% eleve = {"nom": "Alice", "classe": "1NSI", "moyenne": 14.5}
% assert eleve["nom"] == "Alice"
% assert eleve["moyenne"] == 14.5
% assert len(eleve) == 3
% eleve["email"] = "alice@lycee.fr"
% assert eleve["email"] == "alice@lycee.fr"
% assert len(eleve) == 4
% END-VERIFY
\begin{corrige}{1NSI-TC-EX-035}
\begin{enumerate}
  \item Les trois affichages sont :
  \begin{itemize}
    \item \lstinline{Alice} --- valeur associée à la clé \lstinline{"nom"}.
    \item \lstinline{14.5} --- valeur associée à la clé \lstinline{"moyenne"}.
    \item \lstinline{3} --- le dictionnaire contient trois paires clé-valeur.
  \end{itemize}

  \item Les clés sont \lstinline{"nom"}, \lstinline{"classe"} et \lstinline{"moyenne"}. Les valeurs correspondantes sont \lstinline{"Alice"}, \lstinline{"1NSI"} et \lstinline{14.5}.

  \item Pour ajouter une nouvelle paire clé-valeur :
\begin{python}
eleve["email"] = "alice@lycee.fr"
\end{python}
\end{enumerate}
\end{corrige}
